\subsection{Hotel Revenue Management}

\paragraph{Definition and goal.}
Contemporary industry guidance defines hotel revenue management as ``\emph{the process of maximizing revenue and profits through the strategic management of pricing, hotel inventory, and distribution channels}'' (tr.~5)~\cite{amadeus2023}: revenue managers analyse market demand, adjust pricing and inventory levels, and employ various techniques to optimise revenue and profit margins. Its goal is ``\emph{to achieve the right balance between occupancy and rate, ensuring that hotels are operating at maximum capacity while also charging optimal prices}'' (tr.~5)~\cite{amadeus2023} --- the same ADR--occupancy tension that RevPAR captures. The practice originated in the airline industry in the 1980s and was adopted by hospitality because of ``\emph{the perishable nature of hotel rooms and the highly competitive market}'' (tr.~7)~\cite{amadeus2023}.

\paragraph{Foundational aspects.}
Four principles structure the discipline (tr.~6)~\cite{amadeus2023}: \emph{segmentation} --- understanding and targeting customer segments based on their needs, preferences, and willingness to pay; \emph{forecasting} --- predicting future demand for rooms and adjusting prices and inventory accordingly; \emph{optimisation} --- finding the optimal pricing and inventory mix to maximise revenue and profits; and \emph{distribution} --- managing sales channels so that rooms are sold through the most effective channels. Three of the four map one-to-one onto the components of this project: forecasting and optimisation become the model base of Section~3, and segmentation becomes the dimensional backbone of the warehouse; distribution is the axis our dataset cannot support, and it is treated as a documented limitation.

\paragraph{Market segments.}
The industry segments travellers into leisure, business, and group business (tr.~8--10)~\cite{amadeus2023}. Leisure demand is analysed by looking back at which traveller types the hotel has been successful with and forward at booking behaviours; business demand is cultivated proactively through local businesses, trade organisations, and convention bureaus; ``bleisure'' travellers blend business and leisure stays; and group business --- the focus of our allocation problem --- has its own tactics, from dedicated group pricing to partnerships with event planners, and a distinctly longer booking window than individual reservations.

\paragraph{Pricing strategies.}
Pricing optimisation is ``\emph{finding the optimal price for each room type, taking into account market demand, competitor pricing, and guest preferences}'' (tr.~11)~\cite{amadeus2023}. In practice it combines \emph{dynamic pricing} --- adjusting rates in real time based on demand, ``\emph{capitaliz[ing] on high-demand periods and attract[ing] price-sensitive customers during low-demand periods}'' --- with length-of-stay and value-added pricing (tr.~11)~\cite{amadeus2023}. Two monitoring concepts recur throughout the guidance: the \emph{booking window} (``\emph{continuously monitor the booking window across your most profitable segments}'', tr.~12) and \emph{pace} (``\emph{leverage occupancy pace data instead of annualized occupancy metrics}'', tr.~12)~\cite{amadeus2023}, both computed as functions of days before arrival. These are exactly the quantities our forecasting and overbooking models consume.

\paragraph{High-quality data.}
Because forecasting drives every downstream decision, the industry is explicit about what qualifies as good data. Five attributes define the standard (tr.~18)~\cite{amadeus2023}: \emph{on-the-books} (``\emph{data that is a confirmed hotel reservation, not a forecast}''), \emph{forward-looking} (business already booked for future stay dates, not a projection), \emph{sanctioned} (``\emph{extracted in partnership with the provider}''), \emph{segmented} (with ``\emph{significant depth levels of segmentation, market, and traveler attributes}''), and \emph{fresh} (``\emph{refreshed frequently, preferably daily}''). The dataset used in this project satisfies the sanctioned and segmented attributes; the on-the-books, forward-looking, and fresh attributes are addressed by design compromises discussed in Section~5.

\paragraph{Technology: the revenue management system.}
Technology matters for four reasons --- data management, forecasting and pricing, distribution, and automation --- and its most significant embodiment is the revenue management system (RMS): software that ``\emph{helps hoteliers analyze demand patterns, optimize pricing strategies, and maximize revenue by automating the process of analyzing data and adjusting room rates}'', using algorithms that account for ``\emph{seasonality, market demand, competitor pricing, and historical data}'' (tr.~22)~\cite{amadeus2023}. Viewed from industry practice, the system built in this project is a miniature RMS: a warehouse replacing the data engine, a model base replacing the algorithms, an interactive interface replacing the recommendation layer --- with the model management subsystem --- the course-adopted topic, introduced in Section~1.3 --- as the ``S'' that turns analytical models into managed, reproducible software. To avoid losing the revenue gap identified in Section~1.1, revenue management must therefore rest on data and models --- that is, on a decision support system.
